\documentclass[12pt,a4paper]{article}
\usepackage[brazil]{babel}
\usepackage[utf8]{inputenc}
\usepackage[T1]{fontenc}
\usepackage{lmodern}
\usepackage{geometry}
\geometry{top=3cm, bottom=2.5cm, left=3cm, right=2.5cm}
\usepackage{graphicx}
\usepackage{array}
\usepackage{tabularx}
\usepackage{threeparttable}
\usepackage{floatrow}
\usepackage{caption}
\usepackage{hyperref}
\usepackage{listingsutf8}
\usepackage{xcolor}
\usepackage{enumitem}
\usepackage{times}

% === CONFIGURAÇÃO DE CAPTIONS (Força legenda NO TOPO) ===
\newfloat{quadro}{tbhp}{loq}
\floatname{quadro}{Quadro}
\floatsetup[quadro]{capposition=top}
\captionsetup[quadro]{
  position=top,
  justification=centering,
  labelsep=endash,
  width=0.75\linewidth,
  font=small
}

% Comando para fonte padronizada
\newcommand{\oautor}{Wagner Oliveira de Araujo}

% Hyperref
\hypersetup{colorlinks=true, linkcolor=blue, citecolor=blue, urlcolor=blue, filecolor=blue}

% Listings com suporte nativo a UTF-8 e mapeamento de acentos (resolve 100% do erro do TeXnicCenter)
\lstset{
  language=Python,
  extendedchars=true,
  inputencoding=utf8,
  literate=
    {á}{{\'a}}1 {ã}{{\~a}}1 {ç}{{\c{c}}}1 {é}{{\'e}}1 {ê}{{\^e}}1
    {í}{{\'i}}1 {ó}{{\'o}}1 {õ}{{\~o}}1 {ú}{{\'u}}1
    {Á}{{\'A}}1 {Ã}{{\~A}}1 {Ç}{{\c{C}}}1 {É}{{\'E}}1 {Í}{{\'I}}1
    {Ó}{{\'O}}1 {Õ}{{\~O}}1 {Ú}{{\'U}}1,
  basicstyle=\ttfamily\small,
  breaklines=true,
  frame=single,
  backgroundcolor=\color{gray!10},
  keywordstyle=\color{blue},
  commentstyle=\color{gray},
  stringstyle=\color{red!70!black},
  numbers=left,
  numberstyle=\tiny\color{gray},
  stepnumber=1,
  numbersep=5pt,
  showstringspaces=false
}

\begin{document}

\title{\textbf{Projeto Final: Sistema de Regressão Espectral com Bagging e Subagging}\\
\small Levantamento e Implementação em Streamlit/Python}
\author{Wagner Oliveira de Araujo}
\date{IFG --- 2026}
\maketitle
\thispagestyle{empty}

\newpage
\tableofcontents
\newpage
%--------------------------LEVANTAMENTO
\section{Levantamento}

O levantamento de requisitos visou catalogar, classificar e priorizar as funcionalidades do sistema legado (Vaadin/Java) para orientar o desenvolvimento em Python/Streamlit. A análise partiu do mapeamento exaustivo da base de código existente, garantindo rastreabilidade e aderência às necessidades de pesquisadores e alunos de quimiometria.

\subsection{Inventário das Funcionalidades Existentes}

\subsubsection{Funções do Sistema}

Conforme detalhado no Quadro~\ref{quadro:funcoes}, o inventário categoriza as funcionalidades em grupos lógicos (Cálculos, Configuração, Visualização, Organização e Exportação), evidenciando que o sistema original focava estritamente na execução algorítmica, sem mecanismos nativos de exportação ou persistência de resultados.

\begin{quadro}[ht]
\centering
\small
\caption{Inventário das funções do sistema}
\label{quadro:funcoes}
\begin{threeparttable}
\begin{tabularx}{\textwidth}{|l|l|X|} \hline
\textbf{Categoria} & \textbf{Funcionalidade} & \textbf{Descrição Técnica} \\ \hline
Cálculos & RLM Simples (Baseline) & Regressão Linear Múltipla padrão sem técnicas de ensemble \\ \hline
Cálculos & RLM + Bagging & Ensemble com amostragem \textit{com} reposição (bootstrap) \\ \hline
Cálculos & RLM + Subagging + CV & Ensemble com amostragem \textit{sem} reposição + otimização por validação cruzada 5-Fold \\ \hline
Configuração & Parâmetro \textit{m} (Bags) & Define número de amostras/ensembles para o método Bagging \\ \hline
Configuração & Parâmetro \textit{k} (Subagging) & Define fração ou quantidade para amostragem no Subagging \\ \hline
Visualização & Métricas de Desempenho & Exibição de R², RMSE, MAE calculados no conjunto de teste \\ \hline
Visualização & Gráficos Comparativos & \textit{Scatter plots} (previsto $\times$ real), resíduos e comparação entre modelos \\ \hline
Organização & Abas de Navegação & Separação lógica em 3 abas: ``Configuração'', ``Resultados'' e ``Gráficos'' \\ \hline
Exportação & Download de Resultados & Botões de download para métricas em CSV (\texttt{metricas.csv}) e log de execução em TXT (\texttt{relatorio.txt}) \\ \hline
\end{tabularx}
\begin{tablenotes}[flushleft]
\footnotesize
\item Fonte:~\oautor
\end{tablenotes}
\end{threeparttable}
\end{quadro}

\subsubsection{Entradas Necessárias}

Para que os algoritmos sejam executados corretamente, é necessário fornecer matrizes espectrais e vetores de referência em formato texto. Os arquivos e parâmetros exigidos estão listados no Quadro~\ref{quadro:entradas}, que especifica a estrutura dimensional esperada para calibração e teste, além dos limites válidos para os hiperparâmetros \textit{m} e \textit{k}.

\begin{quadro}[ht]
\centering
\small
\caption{Entradas necessárias para execução do sistema}
\label{quadro:entradas}
\begin{threeparttable}
\begin{tabularx}{\textwidth}{|l|l|X|} \hline
\textbf{Tipo} & \textbf{Nome/Parâmetro} & \textbf{Descrição} \\ \hline
Arquivo & \texttt{Xcal.txt} & Matriz de espectros de calibração (variáveis independentes) \\ \hline
Arquivo & \texttt{Ycal.txt} & Vetor de valores de referência para calibração \\ \hline
Arquivo & \texttt{Xteste.txt} & Matriz de espectros para predição/teste \\ \hline
Arquivo & \texttt{Yteste.txt} & Vetor de valores reais para validação do conjunto de teste \\ \hline
Parâmetro & \textit{m} (nº de bags) & Lista de inteiros $\ge 1$ separados por vírgula. Ex.: 10, 20, 30, 40, 50 \\ \hline
Parâmetro & \textit{k} (Subagging) & Lista de decimais em $(0, 1]$ separados por vírgula. Ex.: 0.10, 0.20, 0.30, 0.40, 0.50 \\ \hline
Config. & Validação Cruzada & Fixo: 5-Fold. Parâmetro interno de otimização do modelo \\ \hline
\end{tabularx}
\begin{tablenotes}[flushleft]
\footnotesize
\item Fonte:~\oautor
\end{tablenotes}
\end{threeparttable}
\end{quadro}

\subsubsection{Saídas Esperadas}

Após o processamento dos dados de entrada, o sistema gera indicadores quantitativos e representações visuais para suporte à decisão técnica. As saídas geradas pelo sistema são apresentadas no Quadro~\ref{quadro:saidas}, descrevendo os formatos de exibição adotados para métricas, gráficos comparativos, resíduos e feedbacks de interface.

\begin{quadro}[ht]
\centering
\small
\caption{Saídas esperadas do sistema}
\label{quadro:saidas}
\begin{threeparttable}
\begin{tabularx}{\textwidth}{|l|l|X|} \hline
\textbf{Tipo} & \textbf{Conteúdo} & \textbf{Formato de Exibição} \\ \hline
Numérico & R², RMSE, MAE e Tempo (ms) por modelo & Cards destacados por cor + tabela comparativa interativa \\ \hline
Gráfico & Previsto vs. Real por modelo & Gráfico de linha (Plotly) com séries Real e Predito \\ \hline
Gráfico & Dispersão Real $\times$ Predito (3 modelos) & Scatter plot Plotly com linha ideal $y = x$ \\ \hline
Comparativo & Desempenho entre os três modelos & Tabela \texttt{DataFrame} com RMSE, MAE, R² e Tempo \\ \hline
Exportação & Métricas tabulares & Download \texttt{metricas.csv} via \texttt{st.download\_button} \\ \hline
Exportação & Log completo de execução & Download \texttt{relatorio.txt} via \texttt{st.download\_button} \\ \hline
Feedback & Mensagens de erro/validação & \texttt{st.error()} com descrição orientativa \\ \hline
Status & Indicador de processamento & \texttt{st.progress()} com texto descritivo por etapa \\ \hline
\end{tabularx}
\begin{tablenotes}[flushleft]
\footnotesize
\item Fonte:~\oautor
\end{tablenotes}
\end{threeparttable}
\end{quadro}


Com base no inventário das funcionalidades, entradas e saídas, os requisitos funcionais do sistema foram formalizados e organizados em cinco grupos temáticos (RF01--RF05), cada um subdividido em sub-itens que descrevem comportamentos específicos e verificáveis. Essa formalização, apresentada no Quadro~\ref{quadro:rfs}, é o ponto de partida para a priorização MoSCoW, os testes de validação e a matriz de rastreabilidade.

\begin{quadro}[!ht]
\centering
\small
\caption{Especificação dos requisitos funcionais (RF01--RF05)}
\label{quadro:rfs}
\begin{threeparttable}
\begin{tabularx}{\textwidth}{|l|X|} \hline
\textbf{ID} & \textbf{Descrição do Requisito} \\ \hline
\multicolumn{2}{|l|}{\textbf{RF01 --- Entrada e Validação de Dados}} \\ \hline
RF01.1 & O sistema deve permitir o upload dos arquivos \texttt{Xcal.txt}, \texttt{Ycal.txt}, \texttt{Xteste.txt} e \texttt{Yteste.txt} via \texttt{st.file\_uploader} \\ \hline
RF01.2 & O sistema deve aceitar apenas arquivos com extensão \texttt{.txt}, bloqueando outros formatos no widget de upload \\ \hline
RF01.3 & O sistema deve exibir mensagem de erro descritiva em caso de falha na leitura ou conversão dos arquivos, interrompendo a execução via \texttt{st.stop()} \\ \hline
\multicolumn{2}{|l|}{\textbf{RF02 --- Configuração e Otimização}} \\ \hline
RF02.1 & O sistema deve permitir a entrada de múltiplos valores de \textit{m} (inteiros $\ge 1$) separados por vírgula via campo de texto \\ \hline
RF02.2 & O sistema deve permitir a entrada de múltiplos valores de \textit{k} (decimais em $(0,1]$) separados por vírgula via campo de texto \\ \hline
RF02.3 & O sistema deve executar validação cruzada 5-Fold sobre todas as combinações de \textit{m} e \textit{k} para selecionar automaticamente os hiperparâmetros ótimos \\ \hline
\multicolumn{2}{|l|}{\textbf{RF03 --- Execução dos Modelos}} \\ \hline
RF03.1 & O sistema deve executar a Regressão Linear Múltipla simples via pseudoinversa (\texttt{np.linalg.pinv}) como modelo baseline \\ \hline
RF03.2 & O sistema deve executar RLM com Bagging: amostragem com reposição (\textit{bootstrap}) usando o melhor \textit{m}* obtido por CV \\ \hline
RF03.3 & O sistema deve executar RLM com Subagging: amostragem sem reposição com fração \textit{k}*, usando os melhores \textit{k}* e \textit{m}* obtidos por CV \\ \hline
\multicolumn{2}{|l|}{\textbf{RF04 --- Resultados e Visualização}} \\ \hline
RF04.1 & O sistema deve calcular e exibir RMSE, MAE, R² e tempo de execução em ms para cada modelo via \texttt{calc\_metrics()} e \texttt{time.perf\_counter()} \\ \hline
RF04.2 & O sistema deve exibir cards de métricas por modelo com código de cor (azul/verde/roxo) e tabela comparativa \texttt{DataFrame} \\ \hline
RF04.3 & O sistema deve gerar gráfico de linha (Plotly) com valores reais e preditos por amostra para cada modelo \\ \hline
RF04.4 & O sistema deve gerar gráfico de dispersão Real $\times$ Predito unificado com os três modelos e linha ideal $y = x$ \\ \hline
RF04.5 & O sistema deve disponibilizar download de métricas em \texttt{metricas.csv} e do log em \texttt{relatorio.txt} via \texttt{st.download\_button} \\ \hline
\multicolumn{2}{|l|}{\textbf{RF05 --- Interface e Usabilidade}} \\ \hline
RF05.1 & O sistema deve organizar o fluxo em 3 abas: ``Configuração'' (upload + parâmetros), ``Resultados'' (métricas + log) e ``Gráficos'' \\ \hline
RF05.2 & O sistema deve exibir barra de progresso (\texttt{st.progress}) com texto descritivo por etapa durante o processamento \\ \hline
RF05.3 & O sistema deve persistir os resultados da análise na sessão via \texttt{st.session\_state["result"]}, permitindo navegação entre abas sem reprocessamento \\ \hline
\end{tabularx}
\begin{tablenotes}[flushleft]
\footnotesize
\item Fonte:~\oautor
\end{tablenotes}
\end{threeparttable}
\end{quadro}

\newpage
\subsection{Classificação por Prioridade}

A priorização foi realizada por meio da técnica MoSCoW\footnote{MoSCoW: acrônimo para \textit{Must Have} (obrigatório), \textit{Should Have} (importante), \textit{Could Have} (desejável) e \textit{Won't Have} (fora do escopo atual). Método amplamente adotado em engenharia de software e gestão ágil de projetos para priorização de requisitos.}, separando as funcionalidades em essenciais, desejáveis e extras. Conforme o Quadro~\ref{quadro:essenciais}, os requisitos essenciais (Must Have) garantem o funcionamento mínimo do sistema, abrangendo desde o upload e parse dos arquivos até a exibição de métricas, gráficos e downloads para tomada de decisão.

\subsubsection{Essenciais (Must Have)}

\begin{quadro}[ht]
\centering
\small
\caption{Requisitos essenciais do sistema (Must Have)}
\label{quadro:essenciais}
\begin{threeparttable}
\begin{tabularx}{\textwidth}{|l|X|X|} \hline
\textbf{ID} & \textbf{Requisito} & \textbf{Implementação no app.py} \\ \hline
E01 & Upload dos 4 arquivos (\texttt{.txt}) & \texttt{st.file\_uploader} com \texttt{type=["txt"]}; botão bloqueado enquanto algum arquivo estiver ausente \\ \hline
E02 & Execução da RLM Simples (baseline) & \texttt{perform\_rlm()} via \texttt{np.linalg.pinv(Xcal) @ Ycal} \\ \hline
E03 & Cálculo de métricas (R², RMSE, MAE, Tempo) & \texttt{calc\_metrics()} + \texttt{time.perf\_counter()} para cada modelo \\ \hline
E04 & Execução do Bagging e Subagging com CV & \texttt{perform\_bagging()} + \texttt{find\_optimal\_parameters()} com \textit{grid search} 5-Fold \\ \hline
E05 & Exibição de resultados e downloads & Cards, tabela \texttt{DataFrame}, gráficos Plotly e \texttt{st.download\_button} \\ \hline
\end{tabularx}
\begin{tablenotes}[flushleft]
\footnotesize
\item Fonte:~\oautor
\end{tablenotes}
\end{threeparttable}
\end{quadro}

\newpage
Os requisitos desejáveis não bloqueiam a execução, mas agregam valor à experiência do usuário e à reprodutibilidade. O Quadro~\ref{quadro:desejaveis} lista os itens Should Have indicando quais já estão implementados no \texttt{app.py} e quais permanecem como evolução futura.

\subsubsection{Desejáveis (Should Have)}

\begin{quadro}[ht]
\centering
\small
\caption{Requisitos desejáveis do sistema (Should Have)}
\label{quadro:desejaveis}
\begin{threeparttable}
\begin{tabularx}{\textwidth}{|l|X|X|} \hline
\textbf{ID} & \textbf{Requisito} & \textbf{Situação} \\ \hline
D01 & Gráficos interativos (zoom, tooltip) & \textbf{Implementado} --- Plotly com zoom, pan e tooltip nativos \\ \hline
D02 & Exportação de resultados (CSV, TXT) & \textbf{Implementado} --- \texttt{metricas.csv} e \texttt{relatorio.txt} via \texttt{st.download\_button} \\ \hline
D03 & Tooltips explicativos para \textit{m} e \textit{k} & \textbf{Implementado} --- parâmetro \texttt{help=} nos campos \texttt{st.text\_input()} \\ \hline
D04 & Pré-processamento básico (normalização) & Não implementado --- evolução futura \\ \hline
D05 & Histórico de execuções na sessão & Não implementado --- apenas a última execução é armazenada em \texttt{session\_state} \\ \hline
\end{tabularx}
\begin{tablenotes}[flushleft]
\footnotesize
\item Fonte:~\oautor
\end{tablenotes}
\end{threeparttable}
\end{quadro}

Para fases futuras de evolução do sistema, foram mapeados requisitos que ampliam a escalabilidade e a integração com ecossistemas de dados. Conforme o Quadro~\ref{quadro:extras}, os itens Could Have incluem autenticação, suporte a múltiplos formatos e pipelines automatizadas, alinhando a ferramenta a demandas corporativas e acadêmicas avançadas.

\subsubsection{Extras (Could Have)}

\begin{quadro}[ht]
\centering
\small
\caption{Requisitos extras para fases futuras (Could Have)}
\label{quadro:extras}
\begin{threeparttable}
\begin{tabularx}{\textwidth}{|l|X|X|} \hline
\textbf{ID} & \textbf{Requisito} & \textbf{Cenário de Uso} \\ \hline
X01 & Autenticação e projetos privados & Ambientes corporativos ou pesquisa com dados sensíveis \\ \hline
X02 & Integração com APIs externas & Pipelines automatizadas de quimiometria \\ \hline
X03 & Suporte a formatos (.csv, .xlsx, .mat) & Flexibilidade para ingestão de diferentes fontes \\ \hline
X04 & Otimização automática de hiperparâmetros & Reduz necessidade de conhecimento técnico avançado \\ \hline
X05 & Dashboard comparativo multi-rodada & Análise estatística de estabilidade dos ensembles \\ \hline
\end{tabularx}
\begin{tablenotes}[flushleft]
\footnotesize
\item Fonte:~\oautor
\end{tablenotes}
\end{threeparttable}
\end{quadro}

\subsection{Fluxos de Uso}

\subsubsection{Fluxo Principal}

O fluxo principal de execução segue a sequência implementada no \texttt{app.py}: \textbf{Upload} dos quatro arquivos na aba ``Configuração'' $\rightarrow$ \textbf{Entrada} das listas de \textit{m} e \textit{k} via campos de texto $\rightarrow$ \textbf{Execução} pelo botão ``Executar Análise com Otimização CV'' (habilitado apenas quando os 4 arquivos estão presentes) $\rightarrow$ \textbf{Parse} e leitura dos arquivos $\rightarrow$ \textbf{Grid search} com CV 5-Fold para seleção de hiperparâmetros ótimos $\rightarrow$ \textbf{Avaliação final} dos três modelos no conjunto de teste $\rightarrow$ \textbf{Exibição} de métricas, gráficos e downloads nas abas ``Resultados'' e ``Gráficos''.

\subsubsection{Fluxos Alternativos e Pontos Críticos}

A robustez do sistema depende também do tratamento proativo de falhas e da mitigação de pontos críticos de usabilidade. A matriz de exceções, apresentada no Quadro~\ref{quadro:excecoes}, estabelece comportamentos esperados para entradas inválidas, timeouts e navegação inconsistente, garantindo que o usuário seja orientado a corrigir o fluxo sem perda de contexto.

\begin{quadro}[ht]
\centering
\small
\caption{Tratamento de exceções e usabilidade}
\label{quadro:excecoes}
\begin{threeparttable}
\begin{tabularx}{\textwidth}{|l|X|X|} \hline
\textbf{Cenário} & \textbf{Comportamento do Sistema} & \textbf{Implementação no app.py} \\ \hline
Arquivo não selecionado & Botão de execução desabilitado (\texttt{disabled=not files\_ok}) & Verificação \texttt{all([xcal, ycal, xteste, yteste])} + \texttt{st.info()} orientativo \\ \hline
Arquivo com formato inválido & \texttt{st.error()} com mensagem + \texttt{st.stop()} & Bloco \texttt{try/except} em torno de \texttt{parse\_matrix()} e \texttt{parse\_vector()} \\ \hline
Parâmetros fora do intervalo & \texttt{st.error()} com descrição do valor inválido + \texttt{st.stop()} & Funções \texttt{parse\_int\_list()} e \texttt{parse\_float\_list()} com \texttt{raise ValueError} \\ \hline
Erro na análise & \texttt{st.error()} + \texttt{progress\_bar.empty()} & Bloco \texttt{try/except} em torno de \texttt{run\_full\_analysis()} \\ \hline
Navegação sem executar & Mensagem informativa nas abas Resultados e Gráficos & \texttt{st.info()} condicional a \texttt{"result" not in st.session\_state} \\ \hline\end{tabularx}
\begin{tablenotes}[flushleft]
\footnotesize
\item Fonte:~\oautor
\end{tablenotes}
\end{threeparttable}
\end{quadro}

\subsection{Requisitos Não Funcionais}

Além dos requisitos funcionais, o sistema deve atender a restrições técnicas que garantem qualidade, segurança e desempenho. O Quadro~\ref{quadro:nao-funcionais} especifica os critérios observados na implementação, abrangendo limites de tempo de resposta, clareza das mensagens de erro, isolamento de sessão e comportamento sob múltiplos usuários simultâneos.

\begin{quadro}[ht]
\centering
\small
\caption{Especificação de requisitos não funcionais}
\label{quadro:nao-funcionais}
\begin{threeparttable}
\begin{tabularx}{\textwidth}{|l|X|X|} \hline
\textbf{Categoria} & \textbf{Especificação} & \textbf{Critério de Aceite} \\ \hline
Performance & Tempo de cálculo (RLM) & RLM $< 10\,$s; medido via \texttt{time.perf\_counter()} e exibido em ms \\ \hline
Performance & Tempo de cálculo (Ensemble + CV) & Ensemble $< 2\,$min para os valores padrão de \textit{m} e \textit{k}; barra de progresso indica avanço \\ \hline
Performance & Renderização de gráficos & $< 3\,$s após conclusão do cálculo; renderização delegada ao Plotly no cliente \\ \hline
Usabilidade & Clareza textual e recuperação de erro & Mensagens de \texttt{st.error()} descrevem o valor inválido e orientam a correção \\ \hline
Usabilidade & Tooltips de parâmetros & Parâmetros \textit{m} e \textit{k} possuem \texttt{help=} nos \texttt{st.text\_input()} com exemplos de uso \\ \hline
Segurança & Isolamento de sessão & Dados armazenados em \texttt{st.session\_state} por sessão, isolados pelo Streamlit \\ \hline
Escalabilidade & Crescimento de usuários/dados & Arquitetura \textit{stateless} do Streamlit suporta múltiplas sessões simultâneas independentes \\ \hline
\end{tabularx}
\begin{tablenotes}[flushleft]
\footnotesize
\item Fonte:~\oautor
\end{tablenotes}
\end{threeparttable}
\end{quadro}
%--------------------------LEVANTAMENTO
\newpage
\section{Projeto Lógico}

A estrutura do \texttt{app.py} concentra toda a lógica em um único arquivo, com funções bem definidas e separadas por responsabilidade. O Quadro~\ref{quadro:diagnostico} sintetiza as características arquiteturais identificadas no código, destacando as decisões de implementação adotadas para parsing, cálculo, otimização e apresentação.

\begin{quadro}[ht]
\centering
\small
\caption{Características arquiteturais do app.py}
\label{quadro:diagnostico}
\begin{threeparttable}
\begin{tabularx}{\textwidth}{|l|X|X|} \hline
\textbf{Aspecto} & \textbf{Decisão de Implementação} & \textbf{Função / Trecho no app.py} \\ \hline
Parsing de dados & Leitura linha a linha com substituição de vírgula por ponto & \texttt{parse\_matrix()}, \texttt{parse\_vector()} \\ \hline
Cálculo RLM & Pseudoinversa NumPy: \texttt{pinv(X) @ Y} & \texttt{perform\_rlm()} \\ \hline
Bagging/Subagging & Função unificada com flag \texttt{with\_replacement} e \texttt{seed} configurável & \texttt{perform\_bagging()} \\ \hline
Validação cruzada & 5-Fold com permutação aleatória; retorna média e desvio do RMSE & \texttt{run\_cross\_validation()} \\ \hline
Otimização & Busca exaustiva em grade (\textit{grid search}) sobre listas de \textit{m} e \textit{k} & \texttt{find\_optimal\_parameters()} \\ \hline
Métricas & RMSE, MAE e R² calculados vetorialmente; tempo via \texttt{perf\_counter} & \texttt{calc\_metrics()} \\ \hline
Persistência & Resultado armazenado em \texttt{st.session\_state["result"]} após execução & \texttt{run\_full\_analysis()} \\ \hline
Visualização & Gráficos interativos Plotly: linhas e dispersão com tooltip & \texttt{build\_chart()}, aba Gráficos \\ \hline
Exportação & \texttt{st.download\_button} para \texttt{metricas.csv} e \texttt{relatorio.txt} & Aba Resultados \\ \hline
\end{tabularx}
\begin{tablenotes}[flushleft]
\footnotesize
\item Fonte:~\oautor
\end{tablenotes}
\end{threeparttable}
\end{quadro}

O fluxo lógico implementado segue: \texttt{Upload $\rightarrow$ Parse + validação de formato $\rightarrow$ Configuração de listas de \textit{m} e \textit{k} $\rightarrow$ Grid search com CV 5-Fold $\rightarrow$ Avaliação final no conjunto de teste $\rightarrow$ Exibição de métricas, gráficos e downloads}. Todas as funções de cálculo são independentes da interface Streamlit, recebendo e retornando arrays NumPy, o que as torna testáveis isoladamente.

\newpage
\section{Projeto Físico}

O projeto é implementado em um único arquivo \texttt{app.py}, organizado em blocos funcionais delimitados por comentários. Essa estrutura concentra apresentação e lógica de negócio no mesmo arquivo, com separação clara por funções nomeadas. O Quadro~\ref{quadro:estrutura} descreve cada bloco funcional e sua responsabilidade.

\begin{quadro}[ht]
\centering
\small
\caption{Estrutura funcional do app.py}
\label{quadro:estrutura}
\begin{threeparttable}
\begin{tabularx}{\textwidth}{|l|X|} \hline
\textbf{Bloco / Função} & \textbf{Responsabilidade} \\ \hline
Configuração da página & \texttt{st.set\_page\_config()} + CSS personalizado que replica o visual Vaadin \\ \hline
\texttt{parse\_matrix()} & Leitura de arquivo \texttt{.txt} como matriz NumPy (linhas $\times$ colunas) \\ \hline
\texttt{parse\_vector()} & Leitura de arquivo \texttt{.txt} como vetor NumPy (um valor por linha) \\ \hline
\texttt{perform\_rlm()} & RLM simples via \texttt{np.linalg.pinv(Xcal) @ Ycal} \\ \hline
\texttt{perform\_bagging()} & Bagging e Subagging unificados com flag \texttt{with\_replacement} e \texttt{seed} \\ \hline
\texttt{calc\_metrics()} & Cálculo vetorial de RMSE, MAE e R² \\ \hline
\texttt{run\_cross\_validation()} & Validação cruzada 5-Fold; retorna média e desvio-padrão do RMSE \\ \hline
\texttt{find\_optimal\_parameters()} & \textit{Grid search} sobre listas de \textit{m} e \textit{k}; gera log textual das iterações \\ \hline
\texttt{run\_full\_analysis()} & Orquestra parsing, otimização, avaliação final e montagem do dicionário de resultado \\ \hline
\texttt{parse\_int\_list()} / \texttt{parse\_float\_list()} & Validação e conversão das listas de \textit{m} e \textit{k} digitadas pelo usuário \\ \hline
\texttt{metric\_card\_html()} & Gera HTML dos cards de métricas com cor por modelo \\ \hline
\texttt{build\_chart()} & Cria figura Plotly com séries Real e Predito por modelo \\ \hline
Aba Configuração & Upload dos 4 arquivos + entradas de \textit{m} e \textit{k} + botão de execução \\ \hline
Aba Resultados & Cards, tabela comparativa, log e botões de download \\ \hline
Aba Gráficos & Gráficos de linha por modelo + scatter comparativo \\ \hline
\end{tabularx}
\begin{tablenotes}[flushleft]
\footnotesize
\item Fonte:~\oautor
\end{tablenotes}
\end{threeparttable}
\end{quadro}

O ambiente de execução requer Python 3.10+, Streamlit, NumPy, Pandas e Plotly. Não há dependências externas além dessas bibliotecas, e o sistema não utiliza banco de dados nem sistema de arquivos persistente --- todos os dados ficam em memória durante a sessão.

\newpage
\section{Implementação}

A implementação materializa o projeto lógico e físico em código Python, mantendo fidelidade matemática ao original e ampliando usabilidade e reprodutibilidade. Para garantir a qualidade do código, o Quadro~\ref{quadro:testes} define a estratégia de verificação, estabelecendo metas de cobertura por tipo de teste e critérios de aceite baseados em desempenho e experiência do usuário (SUS $\ge 75$).

\begin{quadro}[ht]
\centering
\small
\caption{Estratégia de testes e critérios de aceite}
\label{quadro:testes}
\begin{threeparttable}
\begin{tabularx}{\textwidth}{|l|X|X|} \hline
\textbf{Tipo de Teste} & \textbf{Escopo} & \textbf{Meta de Cobertura} \\ \hline
Unitário & Funções isoladas (parse, metrics) & $\ge 80\%$ (\texttt{pytest-cov}) \\ \hline
Integração & Fluxo completo: Upload $\rightarrow$ CV $\rightarrow$ Gráficos & 100\% dos caminhos críticos \\ \hline
Performance & Tempo de resposta para datasets típicos & RLM $< 10\,$s; Ensemble $< 2\,$min \\ \hline
Aceite (UAT)\tnote{a} & Validação por pesquisadores/alunos & SUS\tnote{b} $\ge 75$ (Bom/Excelente) \\ \hline
\end{tabularx}
\begin{tablenotes}[flushleft]
\footnotesize
\item Fonte:~\oautor
\item[\textit{a}] UAT (\textit{User Acceptance Testing}): Teste de Aceitação do Usuário --- fase em que usuários reais validam se o sistema atende às suas necessidades antes da entrega final.
\item[\textit{b}] SUS (\textit{System Usability Scale}): questionário padronizado de 10 perguntas criado por John Brooke (1986) para medir a usabilidade percebida de um sistema. O \textit{score} varia de 0 a 100, sendo: $\ge 85$ Excelente, 75--84 Bom, 68--74 Aceitável e $\le 67$ Ruim.
\end{tablenotes}
\end{threeparttable}
\end{quadro}

\newpage
O trecho a seguir ilustra a implementação vetorizada da função de Bagging/Subagging, com \textit{seed} configurável para reprodutibilidade:

\begin{lstlisting}[language=Python]
def perform_bagging(Xcal, Ycal, Xteste, num_bags, with_replacement, fraction, seed=42):
    """RLM + Bagging / Subagging com vetorização e seed configurável."""
    n = Xcal.shape[0]
    sample_size = n if with_replacement else max(1, int(n * fraction))
    rng = np.random.default_rng(seed)
    Ypred = np.zeros(Xteste.shape[0])
    for _ in range(num_bags):
        idx = rng.integers(0, n, size=sample_size) if with_replacement else \
              rng.choice(n, size=sample_size, replace=False)
        beta = np.linalg.pinv(Xcal[idx]) @ Ycal[idx]
        Ypred += (Xteste @ beta) / num_bags
    return Ypred
\end{lstlisting}

\newpage
A rastreabilidade completa é verificada por meio da matriz apresentada no Quadro~\ref{quadro:rtm}, que vincula cada grupo de requisitos às fases clássicas de engenharia de software e aos artefatos de validação correspondentes. Essa matriz assegura que nenhuma funcionalidade levantada seja perdida durante a transição de requisitos para código, permitindo auditoria técnica e controle de escopo em todo o ciclo de desenvolvimento.

\begin{quadro}[!ht]
\centering
\small
\caption{Matriz de rastreabilidade de requisitos (\textit{Requirements Traceability Matrix} - RTM)}
\label{quadro:rtm}
\begin{threeparttable}
\begin{tabularx}{\textwidth}{|l|l|X|} \hline
\textbf{ID Requisito} & \textbf{Fase Clássica} & \textbf{Artefato de Validação} \\ \hline
RF01.1--01.3 & 1, 3, 5 & \texttt{st.file\_uploader}, \texttt{parse\_matrix()}, \texttt{parse\_vector()}, \texttt{st.error()} + \texttt{st.stop()} \\ \hline
RF02.1--02.3 & 1, 2, 4 & \texttt{parse\_int\_list()}, \texttt{parse\_float\_list()}, \texttt{run\_cross\_validation()}, \texttt{find\_optimal\_parameters()} \\ \hline
RF03.1--03.3 & 2, 3, 4 & \texttt{perform\_rlm()}, \texttt{perform\_bagging()}, \texttt{run\_full\_analysis()} \\ \hline
RF04.1--04.5 & 2, 4, 5 & \texttt{calc\_metrics()}, \texttt{metric\_card\_html()}, \texttt{build\_chart()}, \texttt{st.download\_button} \\ \hline
RF05.1--05.3 & 1, 3, 5 & \texttt{st.tabs()}, \texttt{st.progress()}, \texttt{st.session\_state["result"]}, \texttt{st.info()} \\ \hline
\end{tabularx}
\begin{tablenotes}[flushleft]
\footnotesize
\item Fonte:~\oautor
\end{tablenotes}
\end{threeparttable}
\end{quadro}

A implementação resultante atende aos requisitos funcionais e não funcionais levantados, reduz a curva de aprendizado para usuários não especialistas e estabelece base técnica sólida para evoluções futuras (autenticação, pipelines automatizados, suporte a novos formatos). A integração explícita entre texto e quadros, conforme demonstrado ao longo deste documento, garante conformidade com as normas acadêmicas e rastreabilidade técnica completa.

\end{document}